% GENERATED FILE. Edit canonical lessons, then run npm run generate:books.
% canonical-source-hash: fnv1a64:01e5474739d211e2
% canonical-lessons: ES-C13-salir

\chapter{Salir --- Leaving}
\label{ch:spanish-salir}
\hlchaptermodality{72}

\begin{chapteropening}
\textbf{By the end of this chapter:} \emph{I can say that I leave, and use the regular forms for anyone else.}

salgo in the yo slot; sales and sale behave normally.
\end{chapteropening}

\section[salir]{salir — leaving with three singular forms}

\label{lesson:ES-C13-salir}



\begin{quote}
You can retrieve \textbf{pongo, pones, pone}, and you already use \textbf{de} with a place name. Keep that frontier while adding \textbf{salir}, “to leave” or “to go out.”
\end{quote}

\begin{sounds}
The \textbf{g} in \textbf{salgo} is hard. Keep the \textbf{l} clear in every form: \textbf{salgo, sales, sale}.
\end{sounds}

\begin{cousinweb}
\textbf{salir} descends from Latin \textbf{salīre}, “to leap” or “to jump.” The older image of leaping outward developed into the everyday meaning “to go out.” English \textbf{salient} first carried the sense “leaping” or “springing,” and a \textbf{sally} is a sudden movement outward. Those cousins preserve the same motion without requiring a long list of uncertain relatives.
\end{cousinweb}

\begin{grammarlens}[title={Grammar Lens: one more singular set}]
\noindent
\begin{tabularx}{\linewidth}{@{}>{\raggedright\arraybackslash}X>{\raggedright\arraybackslash}X@{}}
\toprule
\textbf{person} & \textbf{present of \textbf{salir}} \\
\midrule
yo & \textbf{salgo} \\
tú & \textbf{sales} \\
él / ella / usted & \textbf{sale} \\
\bottomrule
\end{tabularx}

Say \textbf{salgo · sales · sale}. Only the learned \textbf{yo} form contains \textbf{g}. Compare it with \textbf{pongo} because both forms are now known, not because one form could have predicted the other.

Reuse the known word \textbf{de} and the familiar place name Madrid:

\begin{itemize}
\raggedright
  \item \textbf{Salgo de Madrid.} — “I leave Madrid.”
  \item \textbf{Sales de Madrid.} — “You leave Madrid.”
  \item \textbf{Sale de Madrid.} — “He, she, or formal you leaves Madrid.”
\end{itemize}

Plural persons, dating expressions, house vocabulary, and the noun for “exit” wait until their own lessons.
\end{grammarlens}

\subsection*{Your turn}
Say these aloud:

\begin{itemize}
\raggedright
  \item “salir — salgo, sales, sale”
  \item “Salgo de Madrid.”
  \item “salīre — salient — salir”
\end{itemize}

\begin{tcolorbox}[breakable,colback=teal!4,colframe=teal!35!black,title={Before you move on}]
What does \emph{salir} mean? (“To leave” or “to go out.”) Give the three learned forms. (\emph{Salgo, sales, sale.}) What did Latin \emph{salīre} mean? (“To leap” or “to jump.”) Name one English cousin. (\emph{Salient} or \emph{sally}.) Next: add \textbf{venir} and compare it only with already-learned \textbf{tener}.
\end{tcolorbox}
